\begin{thebibliography}{99}
\bibitem{AKS04}
M.~Agrawal, N.~Kayal, and N.~Saxena,
PRIMES is in P,
\emph{Ann. of Math.} (2) \textbf{160} (2004), no.~2, 781--793.

\bibitem{BG17}
S.~Bettin and S.~M.~Gonek,
The $\theta=\infty$ conjecture implies the Riemann hypothesis,
\emph{Mathematika} \textbf{63} (2017), no.~1, 29--33.
\href{https://doi.org/10.1112/S0025579316000139}
{doi:10.1112/S0025579316000139}.

\bibitem{BIMS17}
V.~Bhargava, G.~Ivanyos, R.~Mittal, and N.~Saxena,
Irreducibility and deterministic $r$-th root finding over finite fields,
in \emph{Proceedings of ISSAC 2017}, ACM, 2017, 37--44.

\bibitem{BR18}
G.~Bhowmik and I.~Z.~Ruzsa,
Average Goldbach and the quasi-Riemann hypothesis,
\emph{Anal. Math.} \textbf{44} (2018), no.~1, 51--56.

\bibitem{BGL}
V.~Blomer, L.~Goldmakher, and B.~Louvel,
$L$-functions with $n$-th-order twists,
\emph{Int. Math. Res. Not.} (2014), no.~7, 1925--1955.
\href{https://doi.org/10.1093/imrn/rns257}
{doi:10.1093/imrn/rns257}.

\bibitem{Dav00}
H.~Davenport,
\emph{Multiplicative Number Theory}, 3rd ed., revised by
H.~L.~Montgomery, Graduate Texts in Mathematics \textbf{74},
Springer, New York, 2000.

\bibitem{DR}
A.~Dunn and M.~Radziwi\l\l,
Bias in cubic Gauss sums: Patterson's conjecture,
\emph{Ann. of Math.} (2) \textbf{200} (2024), no.~3, 967--1057.
\href{https://doi.org/10.4007/annals.2024.200.3.3}
{doi:10.4007/annals.2024.200.3.3}.

\bibitem{EKN20}
A.-S.~Elsenhans, J.~Kl\"uners, and F.~Nicolae,
Imaginary quadratic number fields with class groups of small exponent,
\emph{Acta Arith.} \textbf{193} (2020), no.~3, 217--233.

\bibitem{Ford02}
K.~Ford,
Zero-free regions for the Riemann zeta function,
in \emph{Number Theory for the Millennium, II} (Urbana, IL, 2000),
A K Peters, Natick, MA, 2002, 25--56.
Corrected version: \href{https://arxiv.org/abs/1910.08205}
{arXiv:1910.08205}.

\bibitem{FLN10}
E.~Frenkel, R.~Langlands, and B.~C.~Ng\^o,
Formule des traces et fonctorialit\'e: le d\'ebut d'un programme,
\emph{Ann. Sci. Math. Qu\'ebec} \textbf{34} (2010), no.~2, 199--243.
\href{https://arxiv.org/abs/1003.4578}{arXiv:1003.4578}.

\bibitem{Galbraith12}
S.~D.~Galbraith,
\emph{Mathematics of Public Key Cryptography},
Cambridge University Press, 2012.

\bibitem{GZ23}
P.~Gao and L.~Zhao,
Mean squares of quadratic twists of the M\"obius function,
\emph{J. Number Theory} \textbf{247} (2023), 1--14.
\href{https://doi.org/10.1016/j.jnt.2022.12.007}
{doi:10.1016/j.jnt.2022.12.007}.

\bibitem{GL}
L.~Goldmakher and B.~Louvel,
A quadratic large sieve inequality over number fields,
\emph{Math. Proc. Cambridge Philos. Soc.} \textbf{154} (2013),
no.~2, 193--212.
\href{https://doi.org/10.1017/S0305004112000370}
{doi:10.1017/S0305004112000370}.

\bibitem{Gro13}
T.~H.~Gr\"onwall,
Sur les s\'eries de Dirichlet correspondant \`a des caract\`eres complexes,
\emph{Rend. Circ. Mat. Palermo} \textbf{35} (1913), 145--159.
\href{https://doi.org/10.1007/BF03015596}{doi:10.1007/BF03015596}.

\bibitem{Gru74}
F.~Grube,
Ueber einige Euler'sche S\"atze aus der Theorie der quadratischen Formen,
\emph{Z. Math. Phys.} \textbf{19} (1874), 492--519.

\bibitem{Had96}
J.~Hadamard,
Sur la distribution des z\'eros de la fonction $\zeta(s)$ et ses
cons\'equences arithm\'etiques,
\emph{Bull. Soc. Math. France} \textbf{24} (1896), 199--220.
\href{https://doi.org/10.24033/bsmf.545}{doi:10.24033/bsmf.545}.

\bibitem{Has50}
{\raggedright
H.~Hasse,
\emph{Vorlesungen \"uber Zahlentheorie},
Grundlehren der mathematischen Wissenschaften \textbf{59},
Springer, Berlin, 1950.
\href{https://doi.org/10.1007/978-3-642-52795-1}
{doi:10.1007/978-3-642-52795-1}.
\par}

\bibitem{HB95}
D.~R.~Heath-Brown,
A mean value estimate for real character sums,
\emph{Acta Arith.} \textbf{72} (1995), no.~3, 235--275.

\bibitem{HB00}
D.~R.~Heath-Brown,
Kummer's conjecture for cubic Gauss sums,
\emph{Israel J. Math.} \textbf{120} (2000), 97--124.
\href{https://doi.org/10.1007/s11856-000-1273-y}
{doi:10.1007/s11856-000-1273-y}.

\bibitem{HBP79}
D.~R.~Heath-Brown and S.~J.~Patterson,
The distribution of Kummer sums at prime arguments,
\emph{J. Reine Angew. Math.} \textbf{310} (1979), 111--130.
\href{https://doi.org/10.1515/crll.1979.310.111}
{doi:10.1515/crll.1979.310.111}.

\bibitem{Hec20}
E.~Hecke,
Eine neue Art von Zetafunktionen und ihre Beziehungen zur Verteilung
der Primzahlen, I, II,
\emph{Math. Z.} \textbf{1} (1918), 357--376;
\textbf{6} (1920), 11--51.

\bibitem{IR90}
K.~Ireland and M.~Rosen,
\emph{A Classical Introduction to Modern Number Theory}, 2nd ed.,
Graduate Texts in Mathematics \textbf{84}, Springer, New York, 1990.

\bibitem{OpenAI26b}
OpenAI,
The Quasi-Riemann Hypothesis,
preprint, October 5, 2026.

\bibitem{IK04}
H.~Iwaniec and E.~Kowalski,
\emph{Analytic Number Theory},
American Mathematical Society Colloquium Publications \textbf{53},
American Mathematical Society, Providence, RI, 2004.
\href{https://doi.org/10.1090/coll/053}{doi:10.1090/coll/053}.

\bibitem{Kani11}
E.~Kani,
Idoneal numbers and some generalizations,
\emph{Ann. Sci. Math. Qu\'ebec} \textbf{35} (2011), no.~2, 197--227.

\bibitem{Kor58}
N.~M.~Korobov,
Estimates of trigonometric sums and their applications,
\emph{Uspekhi Mat. Nauk} \textbf{13} (1958), no.~4(82),
185--192 (Russian).
\href{https://www.mathnet.ru/eng/rm7458}{Math-Net.Ru: rm7458}.

\bibitem{KLMS20}
E.~Kowalski, Y.~Lin, Ph.~Michel, and W.~Sawin,
Periodic twists of $\mathrm{GL}_3$-automorphic forms,
\emph{Forum Math. Sigma} \textbf{8} (2020), e15.
\href{https://doi.org/10.1017/fms.2020.7}{doi:10.1017/fms.2020.7}.
\href{https://people.math.ethz.ch/~kowalski/gl3-twists.pdf}{Author version}.

\bibitem{Kub69}
T.~Kubota,
\emph{On automorphic functions and the reciprocity law in a number field},
Lectures in Mathematics, Department of Mathematics, Kyoto University,
No.~2, Kinokuniya Book-Store, Tokyo, 1969.
\href{https://hdl.handle.net/2433/84907}{Kyoto University repository}.

\bibitem{Lan03}
E.~Landau,
Neuer Beweis des Primzahlsatzes und Beweis des Primidealsatzes,
\emph{Math. Ann.} \textbf{56} (1903), 645--670.
\href{https://doi.org/10.1007/BF01444310}{doi:10.1007/BF01444310}.

\bibitem{Lit12}
J.~E.~Littlewood,
Quelques cons\'equences de l'hypoth\`ese que la fonction $\zeta(s)$ de
Riemann n'a pas de z\'eros dans le demi-plan $\operatorname{R}(s)>1/2$,
\emph{C. R. Acad. Sci. Paris} \textbf{154} (1912), 263--266.

\bibitem{Lit28}
J.~E.~Littlewood,
On the class-number of the corpus $P(\sqrt{-k})$,
\emph{Proc. London Math. Soc.} (2) \textbf{27} (1928),
no.~1, 358--372.

\bibitem{LZZ26}
R.~F.~Lu, A.~Zaman, and H.~Zhao,
Numerical computations concerning Landau--Siegel zeros,
preprint, 2026.
\href{https://arxiv.org/abs/2602.03626}{arXiv:2602.03626}.

\bibitem{MM09}
H.~Maier and H.~L.~Montgomery,
The sum of the M\"obius function,
\emph{Bull. London Math. Soc.} \textbf{41} (2009), no.~2, 213--226.
\href{https://doi.org/10.1112/blms/bdn119}
{doi:10.1112/blms/bdn119}.

\bibitem{Mil76}
G.~L.~Miller,
Riemann's hypothesis and tests for primality,
\emph{J. Comput. System Sci.} \textbf{13} (1976), no.~3, 300--317.

\bibitem{MV07}
H.~L.~Montgomery and R.~C.~Vaughan,
\emph{Multiplicative Number Theory I: Classical Theory},
Cambridge Studies in Advanced Mathematics \textbf{97},
Cambridge University Press, 2007.

\bibitem{MS02}
M.~Ram Murty and A.~Sankaranarayanan,
Averages of exponential twists of the Liouville function,
\emph{Forum Math.} \textbf{14} (2002), no.~2, 273--291.

\bibitem{DLMF}
{\raggedright
NIST Digital Library of Mathematical Functions,
\S10.43, equation (10.43.19).
\href{https://dlmf.nist.gov/10.43.E19}{https://dlmf.nist.gov/10.43.E19}.
\par}

\bibitem{OpenAI26}
OpenAI,
The Quasi-Riemann Hypothesis: A Zero-Free Half-Plane $\mathrm{Re}(s)>7/8$,
OpenAI Math Release preprint
\href{https://github.com/openai/math/blob/main/preprints/The-Quasi-Riemann-Hypothesis-September-30-2026/paper.pdf}{OAI:The-Quasi-Riemann-Hypothesis-September-30-2026},
2026.

\bibitem{Page35}
A.~Page,
On the number of primes in an arithmetic progression,
\emph{Proc. London Math. Soc.} (2) \textbf{39} (1935), 116--141.
\href{https://doi.org/10.1112/plms/s2-39.1.116}
{doi:10.1112/plms/s2-39.1.116}.

\bibitem{Pat77}
S.~J.~Patterson,
A cubic analogue of the theta series,
\emph{J. Reine Angew. Math.} \textbf{296} (1977), 125--161.
\href{https://doi.org/10.1515/crll.1977.296.125}
{doi:10.1515/crll.1977.296.125}.

\bibitem{PY19}
I.~Petrow and M.~P.~Young,
A generalized cubic moment and the Petersson formula for newforms,
\emph{Math. Ann.} \textbf{373} (2019), nos.~1--2, 287--353.

\bibitem{PY20}
I.~Petrow and M.~P.~Young,
The Weyl bound for Dirichlet $L$-functions of cube-free conductor,
\emph{Ann. of Math.} (2) \textbf{192} (2020), no.~2, 437--486.
\href{https://doi.org/10.4007/annals.2020.192.2.3}
{doi:10.4007/annals.2020.192.2.3}.

\bibitem{Rie59}
B.~Riemann,
Ueber die Anzahl der Primzahlen unter einer gegebenen Gr\"osse,
\emph{Monatsber. K\"onigl. Preuss. Akad. Wiss. Berlin}
(1859), 671--680.
\href{https://www.maths.tcd.ie/pub/HistMath/People/Riemann/Zeta/}
{Original text and English translation}.

\bibitem{Rod56}
K.~A.~Rodosski\u{\i},
On non-residues and zeros of $L$-functions,
\emph{Izv. Akad. Nauk SSSR Ser. Mat.} \textbf{20} (1956),
no.~3, 303--306 (Russian).

\bibitem{Sie35}
C.~L.~Siegel,
\"Uber die Classenzahl quadratischer Zahlk\"orper,
\emph{Acta Arith.} \textbf{1} (1935), 83--86.
\href{https://doi.org/10.4064/aa-1-1-83-86}
{doi:10.4064/aa-1-1-83-86}.

\bibitem{Tat51}
{\raggedright
T.~Tatuzawa,
On a theorem of Siegel,
\emph{Japan. J. Math.} \textbf{21} (1951), 163--178.
\href{https://doi.org/10.4099/jjm1924.21.0_163}
{doi:10.4099/jjm1924.21.0\_163}.
\par}

\bibitem{Tit30}
E.~C.~Titchmarsh,
A divisor problem,
\emph{Rend. Circ. Mat. Palermo} \textbf{54} (1930), 414--429;
correction, \textbf{57} (1933), 478--479.
\href{https://doi.org/10.1007/BF03021203}{doi:10.1007/BF03021203};
\href{https://doi.org/10.1007/BF03017588}{correction}.

\bibitem{VP96}
C.-J.~de la Vall\'ee Poussin,
Recherches analytiques sur la th\'eorie des nombres premiers.
Premi\`ere partie: La fonction $\zeta(s)$ de Riemann et les
nombres premiers en g\'en\'eral,
\emph{Ann. Soc. Sci. Bruxelles}, deuxi\`eme partie,
\textbf{20} (1896), 183--256.

\bibitem{VP99}
C.-J.~de la Vall\'ee Poussin,
Sur la fonction $\zeta(s)$ de Riemann et le nombre des nombres
premiers inf\'erieurs \`a une limite donn\'ee,
\emph{M\'em. Couronn\'es Autres M\'em. Acad. Roy. Sci.
Lettres Beaux-Arts Belgique}, collection in-$8^\circ$,
\textbf{59} (1899), 1--74.
\href{https://doi.org/10.3406/marb.1899.2449}
{doi:10.3406/marb.1899.2449}.

\bibitem{Vin58}
I.~M.~Vinogradov,
A new estimate of the function $\zeta(1+it)$,
\emph{Izv. Akad. Nauk SSSR Ser. Mat.} \textbf{22} (1958),
no.~2, 161--164 (Russian).
\href{https://www.mathnet.ru/eng/im3962}{Math-Net.Ru: im3962}.

\bibitem{Vin85}
I.~M.~Vinogradov,
\emph{Selected Works}, edited by L.~D.~Faddeev et al.,
Springer, Berlin, 1985.

\bibitem{Wei73}
P.~J.~Weinberger,
Exponents of the class groups of complex quadratic fields,
\emph{Acta Arith.} \textbf{22} (1973), 117--124.
\href{https://doi.org/10.4064/aa-22-2-117-124}{doi:10.4064/aa-22-2-117-124}.

\bibitem{Yos87}
A.~Yoshimoto,
On the cubic theta function,
\emph{Nagoya Math. J.} \textbf{105} (1987), 153--167.

\end{thebibliography}
\end{document}

